\title{Controlling Text-to-Image Diffusion by Orthogonal Finetuning}

\begin{abstract}
Large-scale text-to-image diffusion models demonstrate remarkable proficiency in synthesizing photorealistic visuals from textual descriptions. However, steering or customizing these advanced models to address a variety of downstream tasks remains a pressing research challenge. To address this, we propose a principled adaptation technique called Orthogonal Finetuning~(OFT), designed specifically for finetuning text-to-image diffusion models on such tasks. In contrast to prevailing approaches, OFT rigorously preserves hyperspherical energy, a quantity that reflects the pairwise relationships between neurons mapped onto the unit hypersphere. We demonstrate that maintaining this metric is essential for retaining the semantic image generation capacities of these models. To further ensure robustness during adaptation, we introduce Constrained Orthogonal Finetuning (COFT), which augments OFT by enforcing a radius constraint on the hypersphere. Our focus is on two central finetuning applications: subject-driven generation, aimed at producing subject-centric images in new contexts given a few exemplar images and a prompt; and controllable generation, where the model integrates explicit control signals as part of the generation process. Empirical evaluations indicate that our OFT approach consistently surpasses established methods in both image quality and the rate of convergence.
\end{abstract}

\section{Introduction}

Recent advances in text-to-image diffusion models~\cite{saharia2022photorealistic,ramesh2022hierarchical,rombach2022high} have set new benchmarks for high-fidelity image creation directed by natural language instructions. While these systems excel, the use of text prompts alone often lacks the specificity necessary for nuanced, detailed control over the resulting imagery. In this work, we focus on two pivotal directions within the realm of text-to-image synthesis:

\begin{itemize}[leftmargin=*,nosep]

    \item \textbf{Subject-driven generation}~\cite{ruiz2023dreambooth}: Here, given only a handful of subject images, the challenge is to synthesize new representations of that subject, re-contextualized through additional text prompts.
    \item \textbf{Controllable generation}~\cite{zhang2023adding,mou2023t2i}: In this setting, the model is provided with an auxiliary control input (\emph{e.g.}, canny edge maps, segmentation masks), and must generate outputs conforming to both the provided control and accompanying text prompt.
\end{itemize}

At their core, both tasks involve the challenge of efficiently adapting text-to-image diffusion models via finetuning, while still maintaining the generative strengths acquired during pretraining. The criteria for a robust finetuning approach are: (1) \emph{training efficiency}, entailing reduced numbers of trainable parameters and minimized training epochs; and (2) \emph{generalizability preservation}, which refers to upholding the pretrained model's fidelity and diversity in generation. Standard practices in finetuning involve either marginally updating neuron weights using a low learning rate (\emph{e.g.}, \cite{ruiz2023dreambooth}), or augmenting the model with a minor component that incorporates re-parameterized neuron weights (\emph{e.g.}, \cite{hulora2022,zhang2023adding}). Nonetheless, these straightforward methodologies are not sufficient to assure that the generative capabilities from pretraining are retained. Furthermore, there is a notable gap in systematic understanding with respect to constructing effective finetuning schemes and selecting critical hyperparameters, such as training epochs. One central issue is the absence of a reliable metric for evaluating the retention of the pretrained model’s generative capacity. Prevailing finetuning techniques often operate under the implicit premise that a smaller Euclidean distance between the adapted and original models implies superior preservation of pretraining characteristics, which typically motivates the use of very small learning rates. Although this belief might be valid in certain situations, we contend that solely relying on Euclidean proximity to the pretrained weights is inadequate for comprehensively assessing semantic consistency. Consequently, adopting a more structurally informed metric to quantify the divergence between finetuned and pretrained models could significantly improve both the fidelity to pretraining and the stability of the finetuning process.

Drawing upon empirical findings that hyperspherical similarity effectively captures semantic structure~\cite{liu2017deep,liu2018decoupled,chen2020angular}, we utilize hyperspherical energy~\cite{liu2018learning} as a descriptor for the pairwise relationships between neurons. Specifically, hyperspherical energy is computed as the aggregate hyperspherical similarity (\emph{e.g.}, cosine similarity) among all pairs of neurons within a given layer, thereby reflecting the extent to which neuron activations are uniformly spread across the unit hypersphere~\cite{liu2021learning}. Our central hypothesis is that optimal finetuning should result in minimal alteration to the hyperspherical energy relative to the pretrained baseline. While a straightforward approach would be to introduce a regularization term to hold the hyperspherical energy constant during finetuning, this does not ensure effective minimization of its deviation. Instead, we leverage the invariance property of hyperspherical energy, which states that pairwise hyperspherical similarities are preserved when the same orthogonal transformation is applied to every neuron. With this insight, we introduce Orthogonal Finetuning~(OFT), an approach that adapts large-scale text-to-image diffusion models for downstream tasks while keeping their hyperspherical energy unchanged. The principal mechanism involves learning an orthogonal transformation, shared across the layer, so that pairwise angular relationships between neurons are maintained. OFT may also be interpreted as rotating the coordinate basis within a layer. Taking further into account the correlation between Euclidean proximity to the pretrained model and preservation of original performance, we propose an extension, Constrained Orthogonal Finetuning~(COFT), which enforces that the updated parameters reside on a hypersphere of fixed radius, centered at the pretrained neurons.

The rationale behind employing orthogonal transformation for finetuning neurons is partly inspired by concepts from the 2D Fourier transform. This mathematical tool decomposes an image into its magnitude and phase spectra. The phase spectrum, which encapsulates the angular relationship between the input and the basis, retains most of the image's semantic content. Notably, reconstructing an image using its original phase spectrum and an arbitrary magnitude spectrum can still preserve the core semantics. This observation highlights that manipulating neuron orientations is pivotal for semantically altering the resulting image—a central aim of both subject-driven and controllable generative tasks. Nonetheless, allowing arbitrary modifications to neuron directions introduces excessive flexibility, potentially undermining the generative capabilities acquired during pretraining. To limit this flexibility, we advocate for maintaining the pairwise angles between neurons, motivated by the underlying premise that these angles play a vital role in encoding the neural network’s learned knowledge. Under this framework, it becomes natural to employ a shared orthogonal transformation per layer, ensuring that hyperspherical energy remains consistent.

Additional motivation is drawn from the approach of orthogonal over-parameterized training~\cite{liu2021orthogonal}, which involves training classification models from scratch via orthogonal transformations applied to randomly initialized neural networks. Such randomly initialized networks are characterized by a theoretically low expected hyperspherical energy, and the objective in~\cite{liu2021orthogonal} is to sustain this low energy throughout training, as low hyperspherical energy has been shown to promote better generalization in classification tasks~\cite{liu2018learning,lin2020regularizing}. Their findings indicate that orthogonal transformations suffice to enable training of robust and generalizable classifiers. Unlike this prior work, our investigation centers on finetuning text-to-image diffusion models to enhance controllability and downstream generative capabilities. The distinction between OFT and~\cite{liu2021orthogonal} emerges along two lines. First, whereas the previous method targets reducing hyperspherical energy, our OFT method strives to preserve it at the level of the pretrained model, safeguarding the preexisting model structure during finetuning; reducing hyperspherical energy in diffusion model finetuning risks disrupting critical semantic organization. Second, OFT incorporates mechanisms to limit divergence from the pretrained parameters, giving rise to a constrained variant; such constraints are absent in~\cite{liu2021orthogonal}. The essence of effective finetuning lies in balancing adaptability with preservation of prior knowledge, and we contend that our OFT methodology achieves this balance. Our main contributions are as follows:

\begin{itemize}[leftmargin=*,nosep]

    \item We introduce Orthogonal Finetuning (OFT), a new finetuning strategy aimed at enhancing controllability in text-to-image diffusion models. To bolster stability, we also present a constrained variant that restricts the angular divergence relative to the original pretrained model.
    \item Unlike conventional finetuning approaches, OFT adjusts model parameters while rigorously maintaining the hyperspherical energy—an attribute we empirically determine as central to upholding the generative semantics of the pretrained network.
    \item We deploy OFT on two use cases: subject-driven and controllable generation. Through an extensive set of experiments, we demonstrate that OFT yields clear advantages over existing techniques regarding generation fidelity, convergence rate, and robustness during finetuning. In addition, OFT exhibits greater sample efficiency, achieving strong convergence with significantly fewer training images and epochs.
    \item For the task of controllable generation, we develop a novel control consistency metric to measure controllability. This metric functions by inferring the control signal from generated samples and comparing it directly to the original input signal, thereby providing further empirical support for OFT’s capacity for strong control.
\end{itemize}

\section{Related Work}

\textbf{Text-to-image diffusion models}. The field of text-to-image synthesis has witnessed substantial breakthroughs~\cite{saharia2022photorealistic,ramesh2022hierarchical,rombach2022high,gu2022vector,nichol2022glide}, primarily fueled by the advancement of diffusion-based generative frameworks~\cite{ho2020denoising,song2021score,song2021denoising,dhariwal2021diffusion} coupled with sophisticated vision-language modeling~\cite{su2020vlbert,lu2019vilbert,radford2021learning,wang2021simvlm,li2022blip,li2023blip,alayrac2022flamingo,schuhmann2022laion}. Systems like GLIDE~\cite{nichol2022glide} and Imagen~\cite{saharia2022photorealistic} operate in the pixel domain; GLIDE jointly trains a text encoder and a diffusion model on text-image pairs, whereas Imagen relies on a static pretrained text encoder. Stable Diffusion~\cite{rombach2022high} and DALL-E2~\cite{ramesh2022hierarchical}, in contrast, conduct diffusion modeling within a latent representation space. Stable Diffusion leverages VQ-GAN~\cite{esser2021taming} to derive a learned visual codebook as its latent space, while DALL-E2 builds a unified latent embedding of text and images via CLIP~\cite{radford2021learning}. Beyond diffusion models, research in generative adversarial networks~\cite{reed2016generative,zhang2017stackgan,xu2018attngan,li2019controllable} and autoregressive models~\cite{ramesh2021zero,ding2021cogview,wu2022nuwa,yuscaling} has contributed to progress in text-to-image generation. OFT is designed as a broadly applicable, model-agnostic finetuning method suitable for any text-to-image diffusion architecture.

\textbf{Subject-driven generation}. To address subject alteration, approaches such as~\cite{avrahami2022blended,nichol2022glide} integrate user-provided masks as conditional inputs. Inversion-based techniques~\cite{choi2021ilvr,dhariwal2021diffusion,ramesh2022hierarchical,gal2022image} have been introduced to adjust contextual elements while preserving the main subject. Additionally,~\cite{hertz2022prompt} enables both localized and global edits without mask requirements. These solutions, however, often struggle to fully retain distinctive identity characteristics. Methods like Pivotal Tuning~\cite{roich2022pivotal} refine a generator in the vicinity of an inverted latent code and add regularization to safeguard identity features. Similarly,~\cite{nitzan2022mystyle} constructs a person-specific generative face prior from multiple portraits. \cite{casanova2021instance} generates diverse instance variants, though sometimes at the expense of unique subject details. DreamBooth~\cite{ruiz2023dreambooth} customizes diffusion models by introducing specialized tokens and leveraging a reconstruction as well as a class-prior preservation loss, using only a handful of subject images. OFT is implemented within the DreamBooth paradigm, but, instead of standard finetuning with small learning rates, it updates the model parameters through orthogonal transformations.

\textbf{Controllable generation}. Image-to-image translation is fundamentally a subset of controllable generation, with earlier research predominantly relying on conditional generative adversarial networks~\cite{isola2017image,zhu2017toward,wang2018high,park2019semantic,choi2018stargan}. More recently, diffusion models have also been applied to image-to-image translation tasks~\cite{saharia2022palette,voynov2022sketch,wang2022pretraining}. The advent of ControlNet~\cite{zhang2023adding} marks a notable advance in this area, as it introduces a method to guide a pretrained diffusion model by further adapting it to supplementary control cues, resulting in highly effective controllable generation. Similarly, T2I-Adapter~\cite{mou2023t2i} contemporaneously adopts a strategy of finetuning pretrained diffusion models to enhance control over the generation process. In alignment with the setups in \cite{zhang2023adding,mou2023t2i}, we utilize OFT to refine pretrained diffusion models, which consistently delivers superior controllability while demanding less training data and fewer parameters for finetuning. Notably, OFT achieves this without incurring any extra computational burden during test-time inference.

\textbf{Model finetuning}. The practice of refining large pretrained models for downstream applications has surged in popularity in recent years~\cite{devlin2018bert,brown2020language,he2022masked}. Within the broader scope of finetuning, numerous adaptation approaches (\emph{e.g.}, \cite{hulora2022,houlsby2019parameter,pfeiffer2020adapterfusion}) have been extensively explored, particularly in the domain of natural language processing. Among these, LoRA~\cite{hulora2022} stands as the method most closely related to OFT, postulating that updates to weights during finetuning can be effectively captured via low-rank structures. Distinct from this, OFT employs a multiplicative update of neuron weights through a shared orthogonal transformation at the layer level. This methodology not only preserves the pairwise angular relationships between neurons within a given layer but also ensures greater training stability.

\section{Orthogonal Finetuning}

\subsection{Why Does Orthogonal Transformation Make Sense?}

To motivate the use of orthogonal transformation in the finetuning of text-to-image diffusion models, we break the rationale down into two foundational inquiries: (1) the motivation for adjusting the angles (or directions) of neurons during finetuning, and (2) the justification for leveraging orthogonal transformation specifically for modifying these angles.

Addressing the first question, we draw upon prior empirical findings in \cite{liu2018decoupled,chen2020angular} indicating that the angular difference in features serves as a robust indicator of semantic disparity. SphereNet~\cite{liu2017deep} has demonstrated that constraining all neuron outputs to a unit hypersphere during network training can achieve similar, or even superior, representational capacity and generalization, suggesting that neuron orientation encapsulates the essential information present in the input data. To elucidate the relevance of neuron angular information, we present a simple experimental setup depicted in Figure~\ref{fig:toy_ae}: a standard convolutional autoencoder is trained on flower images. During training, feature maps are generated by applying the conventional inner product (where $z$ denotes an element output from convolutional kernel $\bm{w}$ with sliding window input $\bm{x}$). For evaluation, we generate feature maps in three distinct ways: (a) the inner product as in training, (b) leveraging only magnitude information, and (c) using solely angular information. Results in Figure~\ref{fig:toy_ae} reveal that neuron angles can nearly perfectly reconstruct input images, whereas magnitude information provides little to no reconstruction capability. Notably, the cosine activation (c) is not utilized during training—only the inner product. These findings highlight that neuron directions predominantly encode the semantic content of images, making direction adjustment a promising approach for altering semantic representations.

Regarding the second question, the most straightforward strategy for modifying neuron directions is to simultaneously apply rotation or reflection across all neurons within a layer, leading to the use of orthogonal transformations. While it is possible to apply alternative angular transformations—such as assigning distinct rotations to different neurons—we observe that orthogonal transformations strike an effective balance between regularization and flexibility. Furthermore, it has been established in \cite{liu2021orthogonal} that orthogonal transformations offer ample expressive power for neural network training. To further corroborate our claim, we conduct an experiment focused on the regularization benefits of orthogonality constraints. Adopting the ControlNet~\cite{zhang2023adding} setup for controllable generation, we present our findings in Figure~\ref{fig:ortho}. Our standard OFT method demonstrates consistent performance and achieves high-precision control by the end of training (epoch 20). In contrast, the variant of OFT without the orthogonality constraint fails to produce realistic or controllable outputs, underscoring the critical role of the orthogonality constraint in OFT.

\subsection{General Framework}

Typically, the standard approach to finetuning involves employing gradient descent with a low learning rate to make incremental updates to a model or its specific layers. By choosing a small learning rate, the modifications to the pretrained parameters remain slight, thus effectively enforcing that $\thickmuskip=2mu \medmuskip=2mu \|\bm{M} - \bm{M}^0\|$ (where $\bm{M}$ represents the updated parameters and $\bm{M}^0$ the initial pretrained ones) does not become large. The intent behind this strategy is to update the model while only implicitly discouraging significant deviation from the original pretrained weights. While this restriction is intuitive, it may not be sufficiently strict when tuning large-scale models. To enhance control over parameter adaptation, LoRA supplements this approach by imposing a low-rank requirement on the update, namely $\thickmuskip=2mu \medmuskip=2mu \text{rank}(\bm{M} - \bm{M}^0) = r'$, with $r'$ typically set to a small value. Diverging from LoRA, OFT imposes a constraint based on the similarity of neuron pairs via hyperspherical energy: $\thickmuskip=2mu \medmuskip=2mu \| \text{HE}(\bm{M}) - \text{HE}(\bm{M}^0) \| = 0$, where $\text{HE}(\cdot)$ stands for the hyperspherical energy function evaluated over a weight matrix. 

To clarify this, consider the example of a fully connected layer $\thickmuskip=2mu \medmuskip=2mu \bm{W} = \{\bm{w}_1, \ldots, \bm{w}_n\} \in \mathbb{R}^{d \times n}$ with neurons $\bm{w}_i \in \mathbb{R}^d$ and pretrained weights $\bm{W}^0$. The layer's output $\thickmuskip=2mu \medmuskip=2mu \bm{z}\in\mathbb{R}^n$ is determined as $\thickmuskip=2mu \medmuskip=2mu \bm{z} = \bm{W}^\top \bm{x}$, where the input $\thickmuskip=2mu \medmuskip=2mu \bm{x} \in \mathbb{R}^d$. Within this framework, OFT amounts to minimizing the difference in hyperspherical energy between the finetuned and original parameters:

\begin{equation}\label{eq:oft_principle}
\footnotesize
\min_{\bm{W}} \left\| \text{HE}(\bm{W})-\text{HE}(\bm{W}^0)\right\|~~\Leftrightarrow~~\min_{\bm{W}} \bigg{\|} \sum_{i\neq j} \|\hat{\bm{w}}_i-\hat{\bm{w}}_j\|^{-1}-\sum_{i\neq j} \|\hat{\bm{w}}^0_i-\hat{\bm{w}}^0_j\|^{-1}\bigg{\|}
\end{equation}

Here, each $\thickmuskip=2mu \medmuskip=2mu \hat{\bm{w}}_i = \bm{w}_i/\|\bm{w}_i\|$ is a normalized neuron vector, and the hyperspherical energy for $\bm{W}$ is defined as $\thickmuskip=2mu \medmuskip=2mu \text{HE}(\bm{W}):= \sum_{i\neq j} \|\hat{\bm{w}}_i-\hat{\bm{w}}_j\|^{-1}$. It follows immediately that the global minimum of Eq.~\eqref{eq:oft_principle} is zero; this situation occurs whenever $\bm{W}$ differs from $\bm{W}^0$ only through an orthogonal transformation (rotation or reflection), that is, $\thickmuskip=2mu \medmuskip=2mu \bm{W} = \bm{R}\bm{W}^0$ for some orthogonal $\thickmuskip=2mu \medmuskip=2mu \bm{R} \in \mathbb{R}^{d \times d}$ (with $\det(\bm{R}) = 1$ for rotations or $-1$ for reflections). This principle underlies the OF

\begin{equation}\label{eq:oft_general}
    \footnotesize
    \bm{z}=\bm{W}^\top\bm{x}=(\bm{R}\cdot\bm{W}^0)^\top\bm{x},~~~\text{s.t.}~\bm{R}^\top\bm{R}=\bm{R}\bm{R}^\top=\bm{I}
\end{equation}

Here, $\bm{W}$ represents the weight matrix updated by OFT, while $\bm{I}$ refers to the identity matrix. Figure~\ref{fig:block} presents an overview of the OFT mechanism. In line with LoRA's zero initialization strategy, it is crucial for OFT to commence fine-tuning precisely from the original pretrained parameters $\bm{W}^0$. This objective is met by initializing the orthogonal transformation $\bm{R}$ as the identity, thereby ensuring that the starting weights match the pretrained ones. To maintain the orthogonality constraint on $\bm{R}$ throughout optimization, differentiable orthogonalization approaches, as described in \cite{liu2021orthogonal,lezcano2019cheap}, may be employed. Later, we discuss computationally feasible methods for upholding the orthogonality of $\bm{R}$.

\subsection{Efficient Orthogonal Parameterization}\label{sect:eff_ortho}

Although classical orthogonalization processes such as the differentiable Gram-Schmidt procedure are available, they often entail significant computational overhead~\cite{liu2021orthogonal}. To alleviate this burden, we employ the Cayley transform for orthogonal matrix construction. The method formulates the orthogonal matrix as $\thickmuskip=2mu \medmuskip=2mu \bm{R}=(\bm{I}+\bm{Q})(I-\bm{Q})^{-1}$, where $\bm{Q}$ is a skew-symmetric matrix, i.e., $\thickmuskip=2mu \medmuskip=2mu \bm{Q}=-\bm{Q}^\top$. While this approach is computationally efficient, it restricts the resulting orthogonal matrices to those with determinant $1$, namely, members of the special orthogonal group. Empirically, this limitation does not impair the overall performance. Despite the benefits of Cayley parameterization, the representation of $\bm{R}$ can still lead to an excess of parameters for large dimensionalities $d$. To mitigate this, we advocate structuring $\bm{R}$ as a block-diagonal matrix composed of $r$ individual blocks, resulting in the formulation:

\begin{equation}
    \footnotesize
    \bm{R}=\text{diag}(\bm{R}_1,\bm{R}_2,\cdots,\bm{R}_r)=\begin{bmatrix}
    \bm{R}_1\in \text{O}(\frac{d}{r}) &  &\\
    &\ddots & \\
    & & \bm{R}_r\in \text{O}(\frac{d}{r})
    \end{bmatrix}\in \text{O}(d)
\end{equation}

Here, $\text{O}(d)$ represents the orthogonal group in $d$ dimensions, with $\thickmuskip=2mu \medmuskip=2mu \bm{R}\in\mathbb{R}^{d\times d}$ and $\thickmuskip=2mu \medmuskip=2mu \bm{R}_i\in\mathbb{R}^{d/r\times d/r},\forall i$ denoting orthogonal matrices. In the special case where $\thickmuskip=2mu \medmuskip=2mu r=1$, the block-diagonal orthogonal matrix collapses into a standard, fully unconstrained orthogonal matrix. For a $d\times d$ orthogonal matrix, the total parameter count is $\thickmuskip=2mu \medmuskip=2mu d(d-1)/2$, leading to a parameter complexity of $\mathcal{O}(d^2)$. In contrast, an orthogonal matrix arranged in $r$ blocks has $\thickmuskip=2mu \medmuskip=2mu d(d/r-1)/2$ parameters, giving a complexity scaling as $\mathcal{O}(d^2/r)$. If parameter sharing is enabled by setting all blocks identical (\emph{i.e.}, $\thickmuskip=2mu \medmuskip=2mu \bm{R}_i=\bm{R}_j,\forall i\neq j$), this further reduces parameter growth to $\mathcal{O}(d^2/r^2)$. Importantly, each strategy for economizing parameters maintains the overall orthogonality of the resultant matrix $\bm{R}$, thus ensuring hyperspherical energy is preserved.

We now compare the parameter efficiency of OFT and LoRA. LoRA’s trainable parameter count, with a rank $r'$, is given by $\thickmuskip=2mu \medmuskip=2mu r'(d+n)$. Assuming $r$ and $r'$ scale with the neuron dimension $d$ (\emph{e.g.}, $\thickmuskip=2mu \medmuskip=2mu r=r'=\alpha d$ for a constant $\thickmuskip=2mu \medmuskip=2mu 0<\alpha\leq 1$), LoRA has complexity $\mathcal{O}(d^2+dn)$, while OFT achieves a more compact complexity of $\mathcal{O}(d)$. To illustrate this concretely, consider a weight matrix sized $\thickmuskip=2mu \medmuskip=2mu 128\times 128$; with $\thickmuskip=2mu \medmuskip=2mu r'=8$, LoRA requires $2,048$ trainable parameters, while with $\thickmuskip=2mu \medmuskip=2mu r=8$, OFT uses $960$ parameters, assuming no block sharing.

\subsection{Constrained Orthogonal Finetuning}

OFT’s expressivity can be further reduced by limiting the learned model to remain in close proximity to the pretrained weights. Specifically, COFT enforces the following computation in the forward pass:

\begin{equation}\label{eq:coft_general}
    \footnotesize
    \bm{z}=\bm{W}^\top\bm{x}=(\bm{R}\cdot\bm{W}^0)^\top\bm{x},~~~\text{s.t.}~\bm{R}^\top\bm{R}=\bm{R}\bm{R}^\top=\bm{I},~\norm{\bm{R}-\bm{I}}\leq \epsilon
\end{equation}

where $\bm{R}$ is not only enforced to be orthogonal but also constrained to stay within an $\epsilon$-ball of the identity. While the Cayley parameterization, as described in Section~\ref{sect:eff_ortho}, readily handles the orthogonality constraint, incorporating the $\epsilon$-deviation condition is more challenging within this framework. To better understand the Cayley transform’s behavior under these constraints, we invoke the Neumann series for the approximation $\thickmuskip=2mu \medmuskip=2mu \bm{R}=(\bm{I}+\bm{Q})(I-\bm{Q})^{-1} \approx \bm{I}+2\bm{Q}+\mathcal{O}(\bm{Q}^2)$, assuming the Neumann expansion holds.

since the series converges under the operator norm, it is permissible to relocate the constraint $\thickmuskip=2mu \medmuskip=2mu\|\bm{R}-\bm{I}\|\leq\epsilon$ into the Cayley transform itself. This leads to an equivalent constraint on $\bm{Q}$: $\thickmuskip=2mu \medmuskip=2mu \|\bm{Q}-\bm{0}\|\leq\epsilon'$, where $\bm{0}$ designates the zero matrix, and $\epsilon'$ denotes a distinct tolerance hyperparameter from $\epsilon$. The revised constraint on $\bm{Q}$ can be directly imposed using projected gradient descent techniques. To initialize the orthogonal matrix $\bm{R}$ at the identity, we set $\bm{Q}$ to the zero matrix. COFT incorporates two explicit constraints: one that controls minimal change in hyperspherical energy and another that constrains the distance from the original pretrained model. While the latter is commonly enforced only implicitly in typical finetuning regimes, COFT formalizes it as a direct objective. Although OFT already offers strong empirical results, we find that the explicit deviation restriction in COFT further enhances stability during finetuning. Figure~\ref{fig:eps_coft} illustrates the influence of the hyperparameter $\epsilon$ on COFT's efficacy: larger values of $\epsilon$ yield results akin to those from OFT, while smaller values cause the COFT-updated model to increasingly mimic the initial pretrained text-to-image diffusion architecture.

\subsection{Re-scaled Orthogonal Finetuning}

We introduce a straightforward modification of OFT, wherein a per-neuron scaling factor is also learned. This adaptation arises from the insight that rescaling activations does not affect hyperspherical energy, since normalization removes any magnitude effect. Formally, the revised forward computation is given by: $\thickmuskip=2mu \medmuskip=2mu\bm{z}=(\bm{R}\bm{W}^0\bm{D})^\top\bm{x}$\footnote{\textbf{Errata}: The NeurIPS camera-ready manuscript incorrectly states the forward equation for re-scaled OFT as $\thickmuskip=2mu \medmuskip=2mu\bm{z}=(\bm{D}\bm{R}\bm{W}^0)^\top\bm{x}$. The actual implementation is correct; therefore, results shown in Appendix~\ref{app:rescaled} are unaffected.}, where $\thickmuskip=2mu \medmuskip=2mu \bm{D}=\text{diag}(s_1,\cdots,s_n)\in\mathbb{R}^{n\times n}$ is a learnable diagonal matrix with positive elements $s_1, \dots, s_n$. In contrast to standard OFT—where only $\bm{R}$ serves as a trainable parameter—both $\bm{D}$ and $\bm{R}$ are optimized here. This rescaled OFT variation further extends OFT's adaptability with minimal parameter increase. For the core experiments, we retain the original OFT configuration to emphasize the effect of orthogonal transformations, but empirical evidence suggests the re-scaled variant is generally superior (as detailed in Appendix~\ref{app:rescaled}).

\section{Intriguing Insights and Discussions}

\textbf{OFT is broadly applicable across neural architectures}. In essence, OFT is compatible with any neural network structure. For instance, LoRA is most commonly employed on Transformer attention weights~\cite{hulora2022}; to ensure an equitable comparison with LoRA, our experimental setup restricts OFT's application to the same attention parameters. Beyond dense layers, OFT is also highly suitable for adaptation of convolutional layers. Notably, OFT’s block-diagonal $\bm{R}$ matrices have unique interpretations when applied to convolutions—something not paralleled in LoRA. Assigning the number of blocks equal to the input channel count leads to each block acting exclusively on an individual channel, resembling channel-wise or depth-wise convolution~\cite{chollet2017xception}. Conversely, sharing blocks among all channels results in the same orthogonal transformation being applied to each channel. Initial investigations into convolutional finetuning with OFT are presented in Appendix~\ref{app:convolution}.

\textbf{Connection to LoRA}. LoRA mitigates large deviations in the pretrained weight matrix by introducing a low-rank modification, thus preserving much of the original information. In contrast, OFT enforces the applied transformation on the pretrained weights to be orthogonal and of full rank, thereby ensuring that the core knowledge acquired during pretraining is not compromised. The forward computation for OFT can be re-expressed as $\thickmuskip=2mu \medmuskip=2mu \bm{z}=(\bm{R}\bm{W}^0)^\top\bm{x}=(\bm{W}^0+(\bm{R}-\bm{I})\bm{W}^0)^\top\bm{x}$, where $\thickmuskip=2mu \medmuskip=2mu (\bm{R}-\bm{I})\bm{W}^0$ serves a role similar to LoRA’s low-rank update. Given that $\bm{W}^0$ is generally full-rank, OFT also yields a low-rank adjustment if $\thickmuskip=2mu \medmuskip=2mu \bm{R}-\bm{I}$ is constrained to be low-rank. In parallel with LoRA’s use of a tunable rank parameter $r'$, OFT utilizes a diagonal block parameter $r$, which also lowers the number of learnable parameters. Notably, LoRA and OFT approach parameter efficiency from fundamentally different angles: while LoRA leverages low-rank representations, OFT relies on structural sparsity—specifically, block-diagonal orthogonality—for the same purpose.

\textbf{Why OFT converges faster}. As illustrated in Figure~\ref{fig:toy_ae}, the most direct way to alter model semantics is by modifying neuron orientations, an operation OFT is particularly adept at performing. Furthermore, OFT's optimization can be interpreted as adapting neural weights on a hyperspherical manifold, which produces a smoother optimization surface and thus facilitates improved convergence. This benefit is also supported by empirical observations in \cite{liu2021orthogonal}.

\textbf{Why not minimize hyperspherical energy}. Our method notably differs from \cite{liu2021orthogonal} in that we do not target minimization of hyperspherical energy. In classification tasks, redundancy among neurons is undesirable, but ensuring minimal hyperspherical energy causes neurons to be equidistantly distributed on the hypersphere, which can eliminate important information embedded during pretraining. Thus, this objective is misaligned with the goals of effective finetuning.

\textbf{Balancing flexibility and regularization in finetuning.} Our analysis reveals a fundamental trade-off between model flexibility and regularization. Conventional finetuning offers the highest degree of flexibility; however, it tends to compromise stability and can easily result in model collapse. In contrast, OFT, despite its notable simplicity, effectively mediates this balance by maintaining pairwise angular relationships between neurons. The block-diagonal structure further enforces a stricter regularization constraint on the orthogonal matrix.

\textbf{Inference efficiency preserved.} In contrast to methods like ControlNet, our OFT approach does not incur any extra computational cost during inference for the adapted model. At inference time, the orthogonal transformation represented by the matrix $\bm{R}$ can be fused directly into the original pretrained weight matrix $\bm{W}^0$, resulting in the equivalent matrix $\thickmuskip=2mu \medmuskip=2mu \bm{W} = \bm{R}\bm{W}^0$. As a result, the inference performance remains unchanged compared to the original pretrained model.

\section{Experiments and Results}

\textbf{Experimental protocol.} All experiments utilize the Stable Diffusion v1.5 model~\cite{rombach2022high} as the underlying pretrained text-to-image generator. To ensure an equitable comparison, sample images are randomly selected for each approach. For tasks requiring subject-driven image synthesis, we adopt protocols similar to DreamBooth~\cite{ruiz2023dreambooth}. For tasks involving controllable generation, we generally align with methodologies from ControlNet~\cite{zhang2023adding} and T2I-Adapter~\cite{mou2023t2i}. To fairly evaluate performance relative to LoRA, the OFT or COFT adjustments are confined to those same layers modified by LoRA. Additional implementation details and comprehensive results are available in Appendix~\ref{app:settings}.

\subsection{Subject-driven Generation}

\textbf{Experimental configuration.} We establish DreamBooth~\cite{ruiz2023dreambooth} and LoRA~\cite{hulora2022} as reference methods. All comparative methods employ the loss function described in DreamBooth. For both DreamBooth and LoRA, we follow the implementation details and hyperparameter recommendations outlined in their respective papers. Further experimental outcomes can be found in Appendix~\ref{app:settings}, \ref{app:coft_oft}, \ref{app:more_qualitative}, and \ref{app:failure}.

\textbf{Finetuning stability and convergence}. To begin, we analyze the stability and convergence rate during finetuning for DreamBooth, LoRA, OFT, and COFT, as presented in Figure~\ref{fig:teaser} and Figure~\ref{fig:db_convergence}. It is apparent that COFT and OFT provide stable optimization behavior when adapting the diffusion model. After 400 finetuning steps, effective subject control is evident with DreamBooth and both OFT-based methods, whereas LoRA fails to maintain subject identity. By 2000 iterations, DreamBooth tends to collapse, producing distorted images, and LoRA is unable to generate a yellow shirt, often substituting yellow fur instead. Conversely, OFT and COFT maintain stable and reliable control over the output. These findings highlight the rapid yet robust convergence offered by the OFT approach in subject-conditioned generation tasks. The low sensitivity of OFT and COFT to the total number of finetuning steps is particularly noteworthy, as it reduces the need to finely adjust the iteration count for different subjects. With OFT and COFT, one can select a sufficiently large iteration number without detailed tuning. For COFT, with an appropriately chosen $\epsilon$, both the learning rate and iteration count are straightforward to set.

\textbf{Quantitative comparison}. In line with \cite{ruiz2023dreambooth}, we quantitatively compare the models by assessing subject fidelity (using DINO~\cite{caron2021emerging} and CLIP-I~\cite{radford2021learning}), prompt alignment (CLIP-T~\cite{radford2021learning}), and sample diversity (LPIPS~\cite{zhang2018unreasonable}). Specifically, CLIP-I is defined as the average cosine similarity of CLIP embeddings between real and synthesized images, while DINO applies the same measure with ViT S/16 DINO representations. Prompt fidelity is assessed by the average CLIP embedding similarity between the textual prompt and generated outputs (CLIP-T). Diversity is measured by the average LPIPS cosine distance across samples produced for the same prompt and subject. The outcomes in Table~\ref{table:dreambooth_final} demonstrate that both COFT and OFT achieve substantially higher scores on DINO and CLIP-I metrics compared to DreamBooth and LoRA, and exhibit either superior or comparable results for prompt consistency and diversity. Each approach is tested by running finetuning 30 times with distinct random seeds to reduce the influence of stochasticity. Collectively, these evaluations indicate that OFT not only facilitates more stable and efficient optimization, but also consistently results in improved final model performance.

\textbf{Qualitative comparison}. To facilitate a more perceptible grasp of the advantages conferred by OFT, we present several examples, selected at random, that illustrate subject-driven generation in Figure~\ref{fig:dreambooth_final}. In pursuit of experimental fairness, each example is created using the same finetuned model under each respective method, ensuring that the text prompt remains constant across methods and is not separately optimized for any particular outcome. The model chosen for each method is the one with the highest validation CLIP scores. Upon examining the outcomes in Figure~\ref{fig:dreambooth_final}, it is clear that both OFT and COFT are highly effective at maintaining the semantic features of the subject, whereas LoRA frequently does not retain the identity of the subject (\emph{e.g.}, in the bowl case, LoRA completely fails to maintain subject identity). Furthermore, OFT and COFT demonstrate significantly greater precision in manipulating the generated images via text prompts, a capability where DreamBooth—despite its strengths in preserving subject identity—often falls short in adhering to the provided prompt (\emph{e.g.}, see the first row for the bowl example). This comparative analysis qualitatively evidences that the OFT framework excels in balancing both image controllability and subject retention. In addition, OFT's performance remains robust across a range of iteration numbers, showing insensitivity to this parameter; conversely, neither DreamBooth nor LoRA exhibits similar robustness with larger iteration counts. Supplementary qualitative examples are presented in Appendix~\ref{app:more_qualitative}. Human study findings in Appendix~\ref{app:human} provide further confirmation of OFT's superiority.

\subsection{Controllable Generation}

\textbf{Settings}. As baselines, we adopt ControlNet~\cite{zhang2023adding}, T2I-Adapter~\cite{mou2023t2i}, and LoRA~\cite{hulora2022}. The three principal controllable generation tasks discussed in the main text are: Canny edge to image~(C2I) on COCO~\cite{lin2014microsoft}, segmentation map to image~(S2I) using ADE20K~\cite{zhou2017scene}, and landmark to face~(L2F) on CelebA-HQ~\cite{karras2018progressive,xia2021tedigan}. All approaches undergo finetuning on Stable Diffusion~(SD) v1.5 using these three datasets for 20 epochs. Additional outcomes are available in Appendix~\ref{app:more_qualitative},\ref{app:more_control_task},\ref{app:failure}.

\textbf{Convergence}. We assess the convergence rate of ControlNet, T2I-Adapter, LoRA, and COFT on the L2F task, employing both numerical and visual analyses. For the quantitative measure, we calculate the mean $\ell_2$ distance between the predicted face landmarks and the target control landmarks. In Figure~\ref{fig:face_convergence}, we illustrate how face landmark error evolves across different training epochs for each finetuned model. The data reveals that COFT and OFT both demonstrate considerably quicker convergence in comparison to the baselines. Notably, LoRA requires 20 epochs to match the performance level that OFT reaches by epoch 8. It is important to highlight that OFT and COFT utilize a similar count of trainable parameters as LoRA—substantially less than ControlNet—yet achieve more rapid convergence than alternative approaches. The accelerated learning efficiency of OFT is further corroborated by Figure~\ref{fig:teaser}, where the rightmost example illustrates that OFT remains highly data-efficient: it reaches strong performance using just 5\% of the ADE20K dataset, outperforming ControlNet and LoRA under data constraints. For qualitative comparison, our focus is on OFT, COFT, and ControlNet, as ControlNet exhibits landmark errors comparable to those of our models. The samples in Figure~\ref{fig:face_convergence_qual} indicate that both OFT and COFT not only converge reliably but also progressively align generated face poses to the control landmarks. By contrast, ControlNet’s controllability emerges abruptly after epoch 8, consistent with the sharp convergence effect discussed in \cite{zhang2023adding}. This reliable, smooth convergence characteristic of OFT simplifies practical training and hyperparameter tuning, since OFT exhibits more stable and predictable training behavior.

\textbf{Quantitative comparison}. To assess the effectiveness of controllable generation, we propose a metric termed control consistency, which measures the alignment between the intended input control signal and the control information extracted from the resulting image. For the C2I task, intersection-over-union (IoU) and F1 score are used as evaluation metrics. The S2I task utilizes mean IoU alongside both mean and overall accuracy. For the L2F task, we calculate the mean $\ell_2$ distance between input control landmarks and their predicted counterparts. Further details on these consistency measures are provided in Appendix~\ref{app:settings}. When benchmarking against other approaches, optimal hyperparameter configurations are applied in all cases. As seen in Table~\ref{table:control_gen}, OFT and COFT consistently provide stronger and more precise control compared to alternatives. Our analysis reveals that methods relying on adapters (such as T2I-Adapter and ControlNet) demonstrate slower convergence and ultimately inferior performance. Notably, LoRA surpasses ControlNet in the S2I scenario but falls behind in both C2I and L2F tasks. Overall, the results suggest that LoRA is subject to a relatively low upper limit in performance, even under extensive hyperparameter optimization. In contrast, OFT does not appear to have reached a plateau, as its performance continues to improve with increased parameter count. It is important to note that our proposed quantitative framework for evaluating controllable generation constitutes a key novel contribution, as it reliably measures control fidelity of finetuned models within the limitations of the underlying segmentation or detection models.

\textbf{Qualitative comparison}. We further conduct a qualitative assessment of OFT and COFT against leading approaches such as ControlNet, T2I-Adapter, and LoRA. As depicted by the randomly generated samples in Figure~\ref{fig:control_final}, OFT and COFT not only maintain a high level of image realism and fidelity, but also demonstrate precise control over the generation process. Within the S2I task, it is evident that LoRA is unable to follow the input segmentation map, failing to produce appropriate images, while ControlNet, OFT, and COFT are able to effectively constrain the outputs according to the given control signal. OFT and COFT, when compared to ControlNet, produce visually richer images with more accurate geometries and detailed features, despite employing significantly fewer model parameters. For the C2I task, both OFT and COFT excel in synthesizing images that are semantically consistent with the provided rough Canny edges, outperforming T2I-Adapter and LoRA by a wide margin. Regarding the L2F task, our framework achieves the best results in aligning generated facial images to challenging pose controls. Across all three types of control, the results substantiate that OFT and COFT generate qualitatively superior images relative to existing state-of-the-art methods, thereby validating the effectiveness of our OFT approach for controllable image generation. To supplement these findings, additional qualitative results spanning all three control scenarios can be found in Appendix~\ref{app:control_exp}, and further demonstrations of OFT’s versatility on additional control tasks (such as dense pose to human body, sketch to image, and depth to image) are presented in Appendix~\ref{app:more_control_task}.

\section{Concluding Remarks and Open Problems}

Prompted by the insight that the angular relationships among neurons fundamentally shape visual semantics, we introduce a straightforward yet powerful finetuning approach—orthogonal finetuning—for exerting control over text-to-image diffusion models. Our method specifically addresses two text-to-image tasks: subject-driven generation and controllable generation. Relative to prior techniques, OFT achieves greater control and stability during finetuning while utilizing fewer parameters. Crucially, OFT avoids incurring any extra computational cost at inference time, thereby supporting efficient model deployment.

However, OFT also raises several compelling questions for future exploration. First, the approach ensures orthogonality using Cayley parametrization, which relies on computing a matrix inverse. This requirement can modestly constrain the scalability of OFT. While block diagonal parametrization mitigates some of this issue, developing a faster, differentiable solution for the matrix inversion step remains unresolved. Second, OFT possesses distinctive promise in compositionality: orthogonal matrices generated through distinct OFT finetuning processes can be composed via matrix multiplication, preserving orthogonality. Investigating whether these composed matrices can collectively retain knowledge from all associated downstream tasks is a promising avenue for future research. Finally, the parameter efficiency of OFT currently hinges on the block diagonal design, which may inadvertently introduce biases and restrict expressiveness. Identifying more effective, less biased strategies to enhance parameter efficiency constitutes another important direction for future work.

\end{document}